\section{Introduction}
\label{sec:introduction}

Session types
\cite{DBLP:conf/concur/Honda93,DBLP:conf/parle/TakeuchiHK94,DBLP:conf/esop/HondaVK98}
have become a quasi-standard for defining stateful communication
protocols and reasoning about them. A (binary) session
type specifies a protocol as a set of sequences of typed messages going forth
and back between two participants. Many interesting protocols are
recursive, which means they admit an infinite set of sequences or even
infinite sequences of messages. Such protocols are modeled using equirecursive session
types like \lstinline{IntListS} (session type for sending finitely many integers and then
close) and \lstinline{IntStreamS} (send infinitely many integers) that
are assigned to one end of a communication channel:
\begin{lstlisting}
  type IntListS   = mualpha. oplus{Cons: !Int.alpha, Nil: EndT} -- finite sequences
  type IntStreamS = mualpha. !Int.alpha                    -- infinite sequence
\end{lstlisting}
In session types, the $\oplus\{l_i:S_i\}$ constructor stands for
selection: send a label $l_i$, like
\lstinline|Cons| or \lstinline|Nil|, and continue according to the
chosen branch $S_i$; the $!T.S$ constructor stands for sending a value of
type $T$ and continuing the protocol according to $S$. The other end
of the channel has the \emph{dual type} where sending and receiving ($?T.S$)
is flipped, as well as selection and branching indicated by
$\&\{l_i:S_i\}$: if we receive a label $l_i$, then continue with $S_i$.

Using recursive types directly is somewhat tedious as equivalence of
equirecursive types\footnote{An equirecursive type is considered
  \emph{equal} to its unfoldings.} amounts to comparing infinite regular trees or
reasoning up to unfolding of the recursive types. While equivalence is
decidable in $O(n^2)$ for traditional, tail-recursive session types
(where $n$ corresponds to the size of the type), 
the difficulty increases when considering context-free session types
\cite{DBLP:conf/icfp/ThiemannV16,DBLP:conf/esop/Padovani17,DBLP:journals/toplas/Padovani19,DBLP:journals/corr/abs-2106-06658,
  DBLP:conf/esop/DasDMP21} where the first equivalence test was doubly
exponential\cite{DBLP:conf/tacas/AlmeidaMV20}, though recent work
indicates it might be polynomially decidable in
$O(n^{12})$ \cite{DBLP:journals/corr/abs-2407-04063}.

Prior work proposed parameterized algebraic
protocols \cite{DBLP:journals/pacmpl/MordidoS0V23}, where recursive
protocols are defined analogously to algebraic datatypes (ADTs). In this framework, we
can define \lstinline{IntListP}, the algebraic protocol analogous
to \lstinline{IntListS}, and \lstinline{IntStreamP}, analogous to
\lstinline{IntStreamS}, as follows:
\begin{lstlisting}
  protocol IntListP   = Nil | Cons Int IntListP
  protocol IntStreamP = Next Int IntStreamP
\end{lstlisting}
As with ADTs, the combination of recursive types with sum-of-product
types is ``shrink wrapped'' into a nominal type. In consequence,
deciding equivalence of protocol types takes time linear in the size
of the type, regardless of whether the underlying protocol is
tail-recursive or context free.

The prior work gave an informal outline of an embedding of
algebraic protocols into context-free session types, but it did not offer
a comparison with traditional, tail-recursive session types. The
present work fills this gap as follows.

\subsection{Traditional Session Types to Algebraic Protocols}
\label{sec:trad-sess-types}

The main difficulty in translating a program with traditional session
types (TST) to algebraic session types (AST) lies in managing the unfolding of
equirecursive types. That is, at every point in a TST typing
derivation we may silently apply a type conversion that unrolls a
recursive session type as in $\mu\alpha.S \approx S[\mu\alpha.S /
\alpha]$. Our translation to AST proceeds as follows.
First, we translate \emph{all} types in a TST program into a system of recursive
type equations where all left hand sides are type variables and
equations are flat, i.e., they consist of a type constructor applied
to type variables as in
\lstinline{alpha->beta} or \lstinline{?alpha.beta}. All type variables
occurring on right hand sides must be defined by a single
equation. Defined type variables on the left hand sides may be used
zero or more times on the right hand sides. Here are two equation
systems (one line each)
corresponding to our running examples.
\begin{lstlisting}
IntStreamS  alpha = !beta.alpha                   beta = Int
IntListS    alpha = (+){Cons: beta, Nil: delta}   beta = !gamma.alpha    gamma = Int   delta = End
\end{lstlisting}
Such an equation system can be considered a finite automaton or a
labeled transition system (LTS), where each variable on the right is
reachable by one labeled transition from the left hand side. For
example, in the LTS corresponding to \lstinline{IntListS}, there are
labeled transitions $\alpha \stackrel{Cons}{\longrightarrow} \beta$ and $\alpha \stackrel{Nil}
{\longrightarrow} \delta$ and further transitions for $\beta$,
$\gamma$, and $\delta$. In a second step, we minimize this
automaton. Minimization corresponds to equating bisimilar type
variables, where bisimilarity corresponds to type equivalence of
equirecursive types. Thus, after minimization two types are bisimilar
if and only if they are represented by the \emph{same} type
variable. Consequently, we can eliminate the conversion rule from
typing as two types are now equivalent iff they are represented by the
same type variable.

The final step of the translation reads algebraic protocols from
the equations. To this end, we generate a decision protocol for each
type variable with a select type or branch type (like $\alpha$ for
\lstinline|IntListS|) and use predefined protocols for sequence, empty
protocol, and repetition to translate sending or receiving session
types. For \lstinline|IntListS| and \lstinline|IntStreamS|, our
translation yields
the algebraic protocols \lstinline{IntListSP} and
\lstinline{IntStreamSP}, shown after the definitions of the auxiliary protocols.
\begin{lstlisting}
  protocol ConsNil a b = Cons a | Nil b   -- select/branch: Cons or Nil
  protocol Seq a b = Seq a b              -- first a then b
  protocol Eps = Eps                      -- empty protocol
  protocol Repeat a = Repeat a (Repeat a) -- infinitely often a

  IntListSP   = Repeat (ConsNil (Seq Int Eps) End)
  IntStreamSP = Repeat (Seq Int Eps)
\end{lstlisting}
Interestingly, the resulting algebraic protocols are
\emph{interoperable} with their TST origins. We take that to mean that
the types are bisimilar when considered as LTS. Interoperability results
from an innocuous change that was hinted at, but not executed in the
prior work \cite{DBLP:journals/pacmpl/MordidoS0V23}: single
constructor protocols do not send or receive their constructor. That
is, the protocols \lstinline|Seq|, \lstinline|Eps|, and
\lstinline|Repeat| are tag-free as they do not transmit a tag.
In our example, it even holds that all three representations of the
types are bisimilar.
\begin{lstlisting}
  IntListS   ~ IntListSP   ~ IntListP
  IntStreamS ~ IntStreamSP ~ IntStreamP
\end{lstlisting}

\subsection{Subtyping for Algebraic Protocols}
\label{sec:subtyp-algebr-prot}


We define a subtyping relation on parameterized algebraic
protocols. This subtyping relation is inspired by prior work on
constructor subtyping \cite{DBLP:conf/esop/BartheF99}, where a
% restrictable variants \cite{DBLP:conf/ecoop/MadsenSL23}
subtype of an algebraic data type is defined by a non-empty subset of its
constructors. In our framework, the protocol \lstinline{IntListP} comes in
three variants: \lstinline{IntListP[Nil]}, \lstinline{IntListP[Cons]},
and \lstinline{IntListP[Nil,Cons]} ($=$ \lstinline{IntListP}). They roughly correspond to the
session types \lstinline|oplus{Nil: EndT}|,
\lstinline|mualpha.oplus{Cons: !Int.alpha}|, and
\lstinline{IntListS}. To define the subtyping relation on these
variants, we have to give the protocol a \emph{direction}: the type
\lstinline{!IntListP} interprets the protocol in the forward direction
(corresponding exactly to \lstinline{IntListS}), whereas
\lstinline{?IntListP} interprets the protocol in the reverse direction
(corresponding exactly to the dual of \lstinline{IntListS}). Taking
directionality into account, we obtain the following subtyping judgments:
\begin{lstlisting}
!IntListP <: !IntListP[Nil]                ?IntListP[Nil]  <: ?IntListP
!IntListP <: !IntListP[Cons]               ?IntListP[Cons] <: ?IntListP
\end{lstlisting}
The full subtyping relation corresponds to the reflexive transitive
closure of the relation defined by these judgments.

The subtyping relation resulting from constructor subtyping on algebraic protocols
is still decidable in linear time in the size of the type if comparing
two sets of constructors takes constant 
time\footnote{This assumption is reasonable because the number of
  constructors is usually bounded by a small constant and a set of
  constructors can be represented by a bit set.}. In comparison,
subtyping for tail-recursive session types is decidable in $O(n^2)$ and
subtyping for context-free session types is undecidable
\cite{DBLP:journals/toplas/Padovani19}. That is, 
algebraic protocols with constructor subtyping carve out a subset of
context-free session types with an efficiently decidable subtyping relation.

\subsection{Traditional Session Types with Subtyping to Algebraic Protocols}
\label{sec:trad-sess-types-1}


The constructor subtyping relation is key in extending the mapping from traditional,
tail-recursive session types to algebraic protocols to a setting with
subtyping. The first step in constructing this mapping is just like in
\cref{sec:trad-sess-types}: we represent all types in a program by a system of type
equations. The second step is also the same: we minimize the
equations so that each type variable represents an equivalence class
of bisimilar types.

Stepping back we realize that in our current model, 
a type is a pair $\langle\alpha, \mathcal{E}\rangle$ where
$\mathcal{E}$ is a system of type equations and $\alpha$ is the left
hand side of an equation in $\mathcal{E}$. To stay in this framework,
we define subsupmtion on type equations: Subsumption keeps the
variable $\alpha$; it only adapts $\mathcal{E} <: \mathcal{E}'$ where
the standard subsumption may be applied to any suitable equation. For instance,
here is the subsumption from \lstinline{IntListS} to
\lstinline{IntListS[Nil]} which eliminates the constructor
\lstinline{Cons}: 
\begin{lstlisting}
IntListS         alpha = (+){Cons: beta, Nil: delta}   beta = !gamma.alpha    gamma = Int   delta = End  
IntListS[Nil]    alpha = (+){Nil: delta}            beta = !gamma.alpha    gamma = Int   delta = End
\end{lstlisting}
Importantly, the domains of $\mathcal{E}$ and $\mathcal{E}'$ are the
same, i.e., all left hand sides stay the same and we do not drop
unreachable type variables like $\beta$ and $\gamma$. Getting from the
right hand side of $\alpha$ in the first line to the second line
consists in applying the standard subsumption rule for selection in
functional session types \cite{DBLP:journals/jfp/GayV10}. In the
general case, we may apply subsumption to each equation individually,
taking care whether the left hand side variable occurs in covariant
(like in our example) or in contravariant position.

\subsection{Contributions}
\label{sec:contributions-1}

Prior work stated that algebraic protocol types are nominal and
isorecursive.

OLD-OLD-OLD

This combination of features reduces the complexity of type equivalence from
asymptotically doubly-exponential \cite{DBLP:conf/tacas/AlmeidaMV20,DBLP:conf/esop/DasDMP21} 
to linear. Notably, the expressivity of algebraic protocols 
is on par with that of session types with non-regular recursion
\cite{DBLP:conf/icfp/ThiemannV16,DBLP:conf/esop/DasDMP21}, 
so all protocols expressible in such theories can be adapted 
to algebraic protocols (without paying the price of hosting complex or 
incomplete algorithms in the compiler). Interestingly, types 
are now also iso-associative, thus saving us from defining 
hard-coded associativity laws, due to the basic properties 
of sequential composition 
\cite{DBLP:conf/icfp/ThiemannV16,DBLP:journals/corr/abs-2106-06658}.

\subsection*{Contributions}
\label{sec:contributions}

\begin{itemize}
\item A subtyping relation for algebraic session types which is generated
  by selecting a non-empty subset of protocol constructors.

  This subtyping relation is defined such that its symmetric kernel is
  the equivalence relation in prior work
  \cite{DBLP:journals/pacmpl/MordidoS0V23}. It is decidable in linear
  time and the prior type soundness proof extends to the system with
  subtyping. 

\item The core calculus \algst which is based
  on linear System~F with (non-linear) recursion and session types
  based on algebraic protocols is an extension of prior work. Type
  checking for \algst is specified by bidirectional typing 
  rules that have a direct algorithmic reading.
  We prove type soundness via preservation and progress.
  
\item Traditional, tail-recursive session types \emph{without subtyping} can
  be mapped to \algst. The mapping preserves typing and gives rise to
  a reduction correspondence. 
\item Traditional, tail-recursive session types \emph{with subtyping}
  can be mapped to \algst if all session types in a program can be
  expressed by a system of recursive type equations in a certain
  form. The mapping preserves typing and gives rise to a reduction
  correspondence. 
  
\item Full implementation as a publicly available artifact
  \cite{artifact}, including practical extensions (cf.\ \cref{sec:implementation}).
\end{itemize}

\subsection*{Overview}
\label{sec:overview}

\Cref{sec:algebr-sess-types} contains an informal introduction of algebraic
protocols with examples; we illustrate the gain in
expressivity and modularity with respect to other work in
\cref{sec:param-prot}. \Cref{sec:types} discusses the type structure
of \algst and establishes key results for efficient type checking: correctness of type normalization and
worst case complexity of type equivalence. \Cref{sec:processes} defines
statics and dynamics of expressions 
and processes % . \Cref{sec:metatheory}
followed by a presentation of the metatheoretical results.
% \Cref{sec:cfst-vs-algst} gives an informal outline of
% an embedding of context-free session types to \algst.
\Cref{sec:implementation} discusses the implementation; since all rules are
algorithmic, their appropriate reading leads directly to a type checker and an interpreter.
\Cref{sec:related,sec:conclusion} discuss related work and conclude.
%
Proofs, auxiliary results, further examples, an informal outline of an embedding
of context-free session types to \algst, and further discussion are provided as supplementary material. The
implementation is publicly available \cite{artifact}.

%%% Local Variables:
%%% mode: latex
%%% TeX-master: "main"
%%% End:
